% Options for packages loaded elsewhere
\PassOptionsToPackage{unicode}{hyperref}
\PassOptionsToPackage{hyphens}{url}
\documentclass[
  ukrainian,
]{article}
\usepackage{xcolor}
\usepackage[margin=2.5cm]{geometry}
\usepackage{amsmath,amssymb}
\setcounter{secnumdepth}{-\maxdimen} % remove section numbering
\usepackage{iftex}
\ifPDFTeX
  \usepackage[T1]{fontenc}
  \usepackage[utf8]{inputenc}
  \usepackage{textcomp} % provide euro and other symbols
\else % if luatex or xetex
  \usepackage{unicode-math} % this also loads fontspec
  \defaultfontfeatures{Scale=MatchLowercase}
  \defaultfontfeatures[\rmfamily]{Ligatures=TeX,Scale=1}
\fi
\usepackage{lmodern}
\ifPDFTeX\else
\setmainfont{DejaVu Serif}
\setsansfont{DejaVu Sans}
\setmonofont{DejaVu Sans Mono}
\setmathfont{STIX Two Math}
\fi
\ifPDFTeX\else
  % xetex/luatex font selection
\fi
% Use upquote if available, for straight quotes in verbatim environments
\IfFileExists{upquote.sty}{\usepackage{upquote}}{}
\IfFileExists{microtype.sty}{% use microtype if available
  \usepackage[]{microtype}
  \UseMicrotypeSet[protrusion]{basicmath} % disable protrusion for tt fonts
}{}
\makeatletter
\@ifundefined{KOMAClassName}{% if non-KOMA class
  \IfFileExists{parskip.sty}{%
    \usepackage{parskip}
  }{% else
    \setlength{\parindent}{0pt}
    \setlength{\parskip}{6pt plus 2pt minus 1pt}}
}{% if KOMA class
  \KOMAoptions{parskip=half}}
\makeatother
\usepackage{graphicx}
\graphicspath{{/Users/kalinicenkovitalijmikolajovic/.codex/worktrees/5c15/Script_for_translate_book/Output_Final/}}
\setkeys{Gin}{width=0.82\linewidth,height=0.82\textheight,keepaspectratio}
\makeatletter
\newsavebox\pandoc@box
\newcommand*\pandocbounded[1]{% scales image to fit in text height/width
  \sbox\pandoc@box{#1}%
  \Gscale@div\@tempa{\textheight}{\dimexpr\ht\pandoc@box+\dp\pandoc@box\relax}%
  \Gscale@div\@tempb{\linewidth}{\wd\pandoc@box}%
  \ifdim\@tempb\p@<\@tempa\p@\let\@tempa\@tempb\fi% select the smaller of both
  \ifdim\@tempa\p@<\p@\scalebox{\@tempa}{\usebox\pandoc@box}%
  \else\usebox{\pandoc@box}%
  \fi%
}
% Set default figure placement to htbp
\def\fps@figure{htbp}
\makeatother
\ifLuaTeX
\usepackage[bidi=basic,shorthands=off,provide=*]{babel}
\else
\usepackage[bidi=default,shorthands=off,provide=*]{babel}
\fi
\ifLuaTeX
  \usepackage{selnolig} % disable illegal ligatures
\fi
\setlength{\emergencystretch}{3em} % prevent overfull lines
\providecommand{\tightlist}{%
  \setlength{\itemsep}{0pt}\setlength{\parskip}{0pt}}
\usepackage{bookmark}
\IfFileExists{xurl.sty}{\usepackage{xurl}}{} % add URL line breaks if available
\urlstyle{same}
\hypersetup{
  pdflang={uk},
  hidelinks,
  pdfcreator={LaTeX via pandoc}}

\author{}
\date{}

\begin{document}

\section{ЗАУВАЖЕННЯ}\label{ux437ux430ux443ux432ux430ux436ux435ux43dux43dux44f}

\begin{enumerate}
\def\labelenumi{\arabic{enumi}.}
\tightlist
\item
  Якщо початкова \emph{швидкість} струни дорівнює нулю, тоді розв'язок
  (20.6) має вигляд
\end{enumerate}

\[
u(x, 0) = \sin(3\pi x/L),
\]

І це можна інтерпретувати так. Припустимо, ми розклали початкове
зміщення струни

\[
\frac{\pi}{2} e^{-a+\omega}
\]

в ряд по синусам

\[
u(x, 0) =
\]

Тогда каждый член такого разложения должен совершать колебания

\[
\frac{\pi}{2} e^{-a+\omega}
\]

Якщо ми підсумуємо всі ці коливання, отримаємо розв'язок нашої проблеми.
Наприклад, припустимо, що

Початкове переміщення струни задається функцією
\(f(x) = \sin(\pi x/L) + 0.5 \sin(3\pi x/ L) + 0,25 \sin(5\pi x/L)\) .

Повна відповідь системи на початкову умову такого типу дорівнює сумі
відповідей на кожен терм, тобто

\[
\begin{cases} 2hx, & 0 \le x \le 0.5, \\ 2h(1-x), & 0.5 \le x \le 1. \end{cases}
\]

\begin{enumerate}
\def\labelenumi{\arabic{enumi}.}
\setcounter{enumi}{1}
\tightlist
\item
  Член із числом n у розкладі (20.6) називається ** модою осциляції**
  або ** p-тою гармонікою**. Використовуючи тригонометричні формули, ці
  гармоніки можна записати як
\end{enumerate}

\[
\frac{\pi}{2} e^{-a+\omega}
\]

де \(R_n\) і \(\delta_n\) --- нові довільні константи (амплітуда та
фазовий кут). Ця нова форма представлення n-го моду виявляється більш
корисною для аналізу коливань. Частота коливань od (рад/с) та моду
визначаються відповідно до формули

\pandocbounded{\includegraphics[keepaspectratio]{images/_page_153_Picture_8.jpeg}}

Какие движения будет совершать струна, если ее отпустить? 6. Решите
задачу о затухающих колебаниях струны

де T --- натяг струни, p --- лінійна густина струни. Зверніть увагу, що
частота n-ї гармоніки у n разів більша за основну частоту (n=1). Ця
властивість частоти не характерна для всіх типів коливань. Приємний звук
струни скриму або гітари зумовлений множинністю обертонів основної
частоти. У барабані частоти гармонік високого порядку не є кратними
основної частоти коливань.

\section{ЗАВДАННЯ}\label{ux437ux430ux432ux434ux430ux43dux43dux44f}

\begin{enumerate}
\def\labelenumi{\arabic{enumi}.}
\tightlist
\item
  Знайти розв'язок задачі (20.1), якщо початкові умови задаються
  формулами
\end{enumerate}

\[
\frac{\pi}{2} e^{-a+\omega}
\]

\(u_t(x, 0) = 0.\) Побудувати графіки розв'язків для різних моментів
часу. Чи є це рішення періодичним у часі? Який у нього період?

\begin{enumerate}
\def\labelenumi{\arabic{enumi}.}
\setcounter{enumi}{1}
\tightlist
\item
  Знайти розв'язок проблеми коливань струн (20.1) за початкових умов
  вигляду
\end{enumerate}

\begin{verbatim}
u_t(x, 0) = \sin(3\pi x/L).  Как выглядит график решения в различные моменты времени?
\end{verbatim}

\begin{verbatim}
- 3. Показати це, якщо на рис. 20.2 значення  \lambda \geqslant 0  , тоді розв'язки X(x) T(t) будуть або необмеженими, або ідентичними нулю.
\end{verbatim}

\begin{itemize}
\tightlist
\item
  Знайти розв'язок проблеми коливань струн за початкових умов вигляду
\end{itemize}

\[
\frac{2ax}{(x^2+a^2)^2}
\]

\begin{verbatim}
u_1(x, 0) = (3\pi \alpha/L)\sin(3\pi x/L).  5. Гітарна струна довжиною L = 1 тягнеться за середню точку значенням h (див. рис.) Початкове положення струни можна встановити як
\end{verbatim}

\begin{enumerate}
\def\labelenumi{\arabic{enumi}.}
\setcounter{enumi}{13}
\tightlist
\item
\end{enumerate}

\[
\frac{\cos(ax), |x| > \pi/2a}{0, |x| > \pi/2a}
\]

\begin{enumerate}
\def\labelenumi{\arabic{enumi}.}
\setcounter{enumi}{15}
\tightlist
\item
\end{enumerate}

Представляется ли вам разумным полученное решение? Удовлетворяет ли оно
уравнению с частными производными, граничным и начальным условиям?

\begin{enumerate}
\def\labelenumi{\arabic{enumi}.}
\setcounter{enumi}{6}
\item
  Как вы будете решать неоднородное уравнение с заданными начальными и
  граничными условиям?
\item
\end{enumerate}

\[
\frac{1-|x|, |x| < 1}{0, |x| > 1}
\]

\begin{enumerate}
\def\labelenumi{\arabic{enumi}.}
\setcounter{enumi}{18}
\tightlist
\item
\item
\end{enumerate}

\# ОСЦИЛЯЦІЇ ПУЧКА (РІВНЯННЯ ЧЕТВЕРТОГО ПОРЯДКУ В ЧАСТИННИХ ПОХІДНИХ)

МЕТА ЛЕКЦІЇ: Показати, як проблема осциляції пучка веде до рівняння
четвертого порядку і як ця задача розв'язується у випадку вільно
спокійних кінців. Проводиться порівняння коливань променя з вібраціями
струни скрипки.

Головна різниця між поперечними коливаннями тонкої балки та поперечними
коливаннями струни полягає в тому, що промінь чинить опір згину. Не
заглиблюючись у механіку тонких балок, припустимо, що, враховуючи опір
згинання (замість рівняння хвилі) веде до рівняння четвертого порядку

(21.1)

\[
\frac{\pi}{2} [\delta(\omega+a) + \delta(\omega-a)]
rge  \alpha^2 = K/\rho  ,
\]

K---модуль сдвига (чем больше K, тем жестче балка, тем выше частота
колебаний),

р --- линейная плотность балки (масса/ед. длины).

Оскільки це перший випадок, коли читач стикається з рівнянням, вищим за
другий порядок у цій книзі, буде корисно розв'язати типову задачу
коливань пучка. Пізніше ми ознайомимося з іншими проблемами теорії
пучків.

\section{Свободно опирающаяся
балка}\label{ux441ux432ux43eux431ux43eux434ux43dux43e-ux43eux43fux438ux440ux430ux44eux449ux430ux44fux441ux44f-ux431ux430ux43bux43aux430}

Розглянемо невеликі коливання тонкої балки, кінці якої вільно лежать на
двох опорах. Коли ми кажемо «вільно відпочивати», ми маємо на увазі, що
кінці балки не рухаються, але

Клони пучка на кінцевих точках можна змінювати (кінці пучка фіксуються
за допомогою штифтових пристроїв, рис. 21.1).

Совершенно ясно, что на концах балки должны выполняться граничные
условия

\[
u_{tt} = -\alpha^2 u_{xxxx},
\]

\(u(1, t)\) = 0, але зовсім не очевидно, що на кінцях балки мають бути
виконані ще дві граничні умови:

\[
f(x) = \int_{0}^{\infty} F(\omega) \sin(\omega x) d\omega
\]

\(u_{xx}\)(1, t) = 0. Використовуючи теорію тонких балок, можна
показати, що момент згинання в балці пропорційний значенню \(u_{xx}\) ,
а момент згинання на вільно лежачому кінці має дорівнювати P I C

\[
u(0, t) = 0,
\]

.21.1. Свободно оппрающаяся балка.

Нуль вен. Відповідно, зображення, зображене на рис. 21.1 Осцилюючий
пучок описується такою змішаною задачею (для простоти значення o
вважається рівним одиниці):

(21.2)

\[
0 < x < \infty
\]

\begin{verbatim}
2.  f(\alpha x) 
\end{verbatim}

\[
u_{xx}(0, t) = 0,
\]

\begin{verbatim}
3.  e^{-\alpha x} 
\end{verbatim}

Щоб розв'язати цю задачу, ми використаємо метод розділення змінних. Ми
будемо шукати лише періодичні розв'язки, тобто флуктуації

(21.3)

\[
-\omega^2 F(\omega) + \frac{2}{\pi} \omega f(0)
\]

Зверніть увагу, що вибір розв'язку вигляду (21.3) фактично означає, що
ми взяли константу розділення у методі розділення змінних як від'ємну.

Подставим теперь (21.3) в уравнение колебаний балки и получим уравнение
для \(X\left(x\right)\)

общее решение которого записывается в виде

\[
\frac{1}{\alpha} F\left(\frac{\omega}{\alpha}\right)
\]

Щоб визначити коефіцієнти C, D, E та F, підставляємо загальний розв'язок
у граничні умови

\[
\frac{2\omega}{\pi (\alpha^2 + \omega^2)}
x^{-1/2}
u(x, t) = X(x) [A \sin(\omega t) + B \cos(\omega t)].
X^{(1V)} - \omega^2 X = 0,
X(x) = C\cos\sqrt{\omega}x + D\sin\sqrt{\omega}x + E \cosh\sqrt{\omega}x + F \sinh\sqrt{\omega}x.
\]

Откуда получаем два уравнения

\[
0 < x < \infty \qquad \qquad 0 < \omega < \infty
\]

Решая их, находим

,\\
\(\sin \sqrt{\omega} = 0 \Rightarrow \omega = (n\pi)^2\) ,
\(n = 1, 2, \dots\)

Другими словами, \emph{собственные частоты} свободно опирающейся балки
равны

\[
S_n = \frac{2}{\pi} \int_0^{\pi} f(x) \sin(nx) dx.
\]

а фундаментальные решения и, (удовлетворяют уравнению и граничным
условиям (21.2)) записываются в виде

\[
y = \frac{\pi}{b - a} (x - a).
\]

А теперь, поскольку уравнение и граничные условия линейны и однородны,
можно утверждать, что сумма

(21.4)

\[
f(x) = \sum_{n=1}^{\infty} S_n \sin(nx).
\]

также удовлетворяет уравнению и граничным условиям. Единственное, что
нам осталось---это выбрать коэффициенты \(a_n\) и \(b_n\) так, чтобы
удовлетворить начальным условиям. Подставляя (21.4) в начальные условия,
получаем

(21.5)

\[
S_n = \frac{2}{\pi} \int_{0}^{\pi} f(x) \sin(nx) dx
n = 1, 2, ...
\]

Воспользовавшись ортогональностью семейства функции (sin (nnx)) на
отреже {[}0, 1{]}, находим формулы для коэффициентов

(21.6)

\(\frac{\pi}{2}\)e\^{}\{ax\}\[
\] Таким образом, решение задачи имеет вид (21.4), а коэффициенты
\(a_n\) и \(b_n\) вычисляются по формулам (21.6)

Чтобы помочь читателю разобраться в этой задаче, приведем простой
пример.

\end{document}
